\subsection{PROPERTIES OF GROUPS ACTING PROPERLY}

PROPOSITION 3. If a topological group G operates properly on a topological space X, then the orbit space X/G is Hausdorff. If also G is Hausdorff, then X is Hausdorff.

Let \(C \subset X \times X\) be the graph of the equivalence relation \(R\) defined by \(G\) on \(X\) ; then \(C\) is the image of \(G \times X\) under the mapping \(\theta : (s, x) \to (x, s.x)\) . Since \(\theta\) is proper, \(C\) is closed in \(X \times X\) (Chapter I, § 10, no. 1, Proposition 1). Since the relation \(R\) is open (§ 2, no. 4, Lemma 2), it follows that \(X/G\) is Hausdorff (Chapter I, § 8, no. 3, Proposition 8).

Now suppose that G is Hausdorff. Then the mapping \(x \to (e, x)\) of X into \(G \times X\) is a homeomorphism of X onto a closed subspace of \(G \times X\) , and is therefore proper (Chapter I, § 10, no. 1, Proposition 2). If we compose this mapping with the mapping \((s, x) \to (x, s.x)\) of \(G \times X\) into \(X \times X\) , which by hypothesis is proper, we obtain a proper mapping of X into \(X \times X\) , namely the diagonal mapping \(x \to (x, x)\) . Hence the diagonal \(\Delta\) of \(X \times X\) is closed in X, and therefore X is Hausdorff.

PROPOSITION 4. Let G be a topological group operating properly on a topological space X, and let x be a point of X. Let G.x denote the orbit of x, and let \(K_{x}\) denote the stabilizer of x. Then:

\begin{enumerate}
\def\labelenumi{\alph{enumi})}
\item
  The mapping \(s \to s.x\) of \(G\) into \(X\) is proper.
\item
  \(\mathbf{K}_x\) is quasi-compact.
\item
  The canonical mapping of \(\mathbf{G} / \mathbf{K}_x\) onto \(\mathbf{G}.x\) is a homeomorphism.
\item
  The orbit G.x is closed in X.
\end{enumerate}

The inverse image of \(\{x\} \times X\) under \(\theta : (s, x) \to (x, s.x)\) is \(G \times \{x\}\) ; hence, by Proposition 3 of Chapter I, § 10, no. 1, the restriction of \(\theta\) to \(G \times \{x\}\) is a proper mapping of \(G \times \{x\}\) into \(\{x\} \times X\) , whence a) follows. Since \(K_x\) is the inverse image of \(x\) under \(s \to s.x\) , b) follows from Chapter I, § 10, no. 2, Theorem 1. c) and d) are consequences of a), by virtue of Chapter I, § 10, no. 1, Propositions 2 and 5 b).

Remark. Proposition 4 shows that if a topological group G operates properly on a homogeneous space G/H, then the subgroup H is quasi-compact (and therefore compact if G is Hausdorff). It can be shown that this is also a sufficient condition for G to operate properly on G/H (Exercise 3).

PROPOSITION 5. Let G (resp. G') be a topological group operating continuously on a topological space X (resp. X'). Let φ be a continuous homomorphism of G into \(\mathbf{G}'\) and let \(\psi\) be a continuous mapping of X into \(\mathbf{X}'\) compatible with \(\varphi\) ( \(\S 2\) , no. 4). Then:

\begin{enumerate}
\def\labelenumi{(\roman{enumi})}
\item
  If \(\varphi\) is surjective and \(\psi\) is surjective and proper, and if \(\mathbf{G}\) operates properly on \(\mathbf{X}\) , then \(\mathbf{G}'\) operates properly on \(\mathbf{X}'\) .
\item
  If \(\varphi\) is proper, if \(\mathbf{G}'\) operates properly on \(\mathbf{X}'\) and if \(\mathbf{X}\) is Hausdorff, then \(\mathbf{G}\) operates properly on \(\mathbf{X}\) .
\end{enumerate}

To prove (i), consider the commutative diagram

\[
\begin{array}{c c c c} \mathbf {G} & \times \mathbf {X} \xrightarrow {\theta} \mathbf {X} & \times \mathbf {X} \\ \alpha & \downarrow & & \downarrow \beta \\ \mathbf {G} ^ {\prime} & \times \mathbf {X} ^ {\prime} \xrightarrow {\theta^ {\prime}} \mathbf {X} ^ {\prime} & \times \mathbf {X} ^ {\prime} \end{array}
\]

where \(\alpha = \varphi \times \psi\) and \(\beta = \psi \times \psi\) . By hypothesis, \(\theta\) is proper; so is \(\beta\) {[}Chapter I, § 10, no. 1, Proposition 4a){]}; hence \(\beta \circ \theta = \theta' \circ \alpha\) is proper {[}Chapter I, § 10, no. 1, Proposition 5a){]}. Since \(\alpha\) is surjective it follows that \(\theta'\) is proper {[}Chapter I, § 10, no. 1, Proposition 5b){]}.

To prove (ii) consider an ultrafilter \(\mathfrak{U}\) on \(\mathbf{G} \times \mathbf{X}\) , such that the mappings

\[
(s, x) \rightarrow s. x \quad \text { and } \quad (s, x) \rightarrow x
\]

converge with respect to \(\mathfrak{U}\) to \(y_0\) and \(x_0\) respectively. It follows that \((s, x) \to \varphi(s). \psi(x)\) and \((s, x) \to \psi(x)\) converge with respect to \(\mathfrak{U}\) . Since \(\mathbf{G}'\) operates properly on \(\mathbf{X}'\) , this implies (no. 1) that \((s, x) \to \varphi(s)\) converges with respect to \(\mathfrak{U}\) to a point \(s_0' \in \mathbf{G}'\) . Since \(\varphi\) is proper, we deduce (Chapter I, § 10, no. 2, Theorem 1) that \((s, x) \to s\) converges with respect to \(\mathfrak{U}\) to a point \(s_0 \in \mathbf{G}\) . The uniqueness of the limit in \(\mathbf{X}\) then shows that \(y_0 = s_0 x_0\) , and hence that \(\mathbf{G}\) operates properly on \(\mathbf{X}\) (no. 1).

